\documentclass[11pt,reqno]{amsart}
\usepackage[margin=1.15in]{geometry}
\usepackage{amsmath,amssymb,amsthm,mathtools}
\usepackage{booktabs}
\usepackage[colorlinks=true,linkcolor=blue,citecolor=blue,urlcolor=blue]{hyperref}
\usepackage{microtype}

\theoremstyle{plain}
\newtheorem{theorem}{Theorem}
\newtheorem{proposition}[theorem]{Proposition}
\newtheorem{corollary}[theorem]{Corollary}
\newtheorem{lemma}[theorem]{Lemma}
\newtheorem{conjecture}[theorem]{Conjecture}
\theoremstyle{definition}
\newtheorem{definition}[theorem]{Definition}
\newtheorem{remark}[theorem]{Remark}
\newtheorem{observation}[theorem]{Observation}

\DeclareMathOperator{\sqf}{sqf}
\newcommand{\Art}{\mathrm{A}}
\newcommand{\Z}{\mathbb{Z}}
\newcommand{\Q}{\mathbb{Q}}

\begin{document}

\title[Mixed correlations of Artin primes]
{Artin status across consecutive primes and distinct bases:\\
mixed exclusion laws and the completion of a $2 \times 2$ programme}

\author{Josh Bald}
\address{Ontario, Canada}
\email{jpbald93@gmail.com}

\subjclass[2020]{Primary 11A07; Secondary 11N05, 11N13, 11Y16, 11Y60}
\keywords{Artin's conjecture, primitive roots, consecutive primes, quadratic
characters, exclusion laws, Lemke Oliver--Soundararajan bias}

\date{\today}

\begin{abstract}
Say a prime $p$ is Artin base $a$ if $a$ is a primitive root modulo $p$. This
paper completes a two-axis programme on the joint statistics of Artin status:
having previously treated one base at consecutive primes, and several bases at
one prime, we now measure the \emph{mixed} correlations
$\delta(a \to b) = P\bigl(p_{n+1} \in \Art_b \mid p_n \in \Art_a\bigr) -
P\bigl(p_{n+1} \in \Art_b \mid p_n \notin \Art_a\bigr)$ for all $144$ ordered
pairs from twelve bases, over all $50{,}847{,}531$ consecutive prime pairs
below $10^9$. On the deterministic side we classify the \emph{mixed exclusion
classes}: gap classes $g$ modulo $\mathrm{lcm}(f_a, f_b)$ at which no
admissible residue permits $\chi_a(p) = \chi_b(p+g) = -1$, so that $p_n$
Artin base $a$ and $p_{n+1}$ Artin base $b$ can never co-occur. Genuinely
mixed instances exist --- for the ordered pairs $(2, 6)$ and $(6, 2)$ the
classes $g \equiv 10, 14 \pmod{24}$ are excluded although neither base
excludes them alone --- and the classification is verified on
$104{,}972{,}810$ pairs with zero exceptions and confirmed complete by a
data-driven converse scan. On the statistical side, the off-diagonal
correlations are small but overwhelmingly significant ($124$ of $132$ with
$|z| > 3$), negative for $101$ of $132$ ordered pairs (mean $-0.0035$), an
order of magnitude weaker than the diagonal ones --- and \emph{asymmetric}:
$\delta(a \to b) \neq \delta(b \to a)$, with differences up to $0.05$,
reflecting the asymmetric conditioning through two distinct integers
$p_n - 1$ and $p_{n+1} - 1$. The strongest organising variable for
$|\delta(a \to b)|$ is the entanglement of the base pair, followed by the
conductor of the \emph{second} base. As a byproduct of extending the base set
we exhibit base $15$ as the first known base whose consecutive-prime Artin
correlation is globally \emph{positive} ($\delta = +0.0122$, $z = 87$): its
three exclusion classes are outweighed by strongly coupled classes, resolving
in the negative the natural conjecture that $\delta(a) < 0$ for every base.
\end{abstract}

\maketitle

\section{Introduction}
\label{sec:intro}

Let $\Art_a$ denote the set of primes for which the integer $a$ is a
primitive root. Artin's conjecture describes the density of $\Art_a$; Hooley
proved it under GRH~\cite{Hooley1967}. This is the fourth paper in a series
studying the \emph{joint} statistics of Artin status computationally at the
scale $x = 10^9$, and it completes a $2 \times 2$ programme:

\begin{center}
\begin{tabular}{lll}
\toprule
 & same base & different bases \\
\midrule
same prime & (trivial) & paper III~\cite{BaldIII} \\
consecutive primes & papers I, II~\cite{BaldI, BaldII} & \textbf{this paper} \\
\bottomrule
\end{tabular}
\end{center}

Papers I and II found that Artin status of a fixed base $a$ at consecutive
primes is anticorrelated, identified deterministic \emph{exclusion laws} ---
gap classes at which two consecutive primes can never both be Artin base $a$
--- and classified them via two mechanisms: reversal of the quadratic
character $\chi_a$, and inadmissibility of the residues that would be needed
for a doubly-negative configuration. Paper III found that Artin statuses of
\emph{different} bases at the \emph{same} prime attract (the shared
factorisation of $p-1$ is a common cause), except on multiplicatively
dependent pairs, where quadratic entanglement produces repulsion and, inside
multiplicative triples, a deterministic three-way exclusion.

The remaining axis --- base $a$ at $p_n$ against base $b$ at $p_{n+1}$ ---
mixes the two settings: the conditioning passes through two \emph{distinct}
integers $p_n - 1$ and $p_{n+1} - 1$ (as in papers I--II), but the two events
involve \emph{distinct} quadratic characters $\chi_a, \chi_b$ (as in paper
III). Both the deterministic and the statistical questions have new content
here.

\subsection*{Deterministic part}
Call a residue class $g$ modulo $L = \mathrm{lcm}(f_a, f_b)$ (where $f_a$ is
the conductor of $\chi_a$) a \emph{mixed exclusion class} for the ordered pair
$(a, b)$ if there is no residue $r$ with $r$ and $r + g$ admissible
(coprime to $L$) and $\chi_a(r) = \chi_b(r + g) = -1$. Since the Artin
property forces the character value $-1$, primes $p_n \equiv r$,
$p_{n+1} = p_n + g'$ with $g' \equiv g \pmod L$ can then never satisfy both
``$a$ is a primitive root of $p_n$'' and ``$b$ is a primitive root of
$p_{n+1}$''. Diagonal instances ($a = b$) are the exclusion laws of papers
I--II. The classification here (Section~\ref{sec:theory}) yields genuinely
mixed instances; the simplest is:

\begin{theorem}[see Section~\ref{sec:theory}]
\label{thm:mixedintro}
For the ordered pairs $(2, 6)$ and $(6, 2)$, the gap classes
$g \equiv 10, 14 \pmod{24}$ are mixed exclusion classes: no consecutive
primes $p_n, p_{n+1} = p_n + g$ with $g \equiv 10, 14 \pmod{24}$ have $2$ a
primitive root of one and $6$ a primitive root of the other. Neither base
excludes these classes by itself.
\end{theorem}

At $x = 10^9$ the predicted mixed exclusion classes across all $144$ ordered
pairs contain $104{,}972{,}810$ consecutive prime pairs and exactly zero
co-occurrences, and a data-driven converse scan (every class of every pair
with at least $5 \times 10^5$ pairs and a zero count) finds nothing the
classification misses.

\subsection*{Statistical part}
Define the mixed correlation
\[
  \delta(a \to b) \;=\;
  P\bigl(p_{n+1} \in \Art_b \mid p_n \in \Art_a\bigr) -
  P\bigl(p_{n+1} \in \Art_b \mid p_n \notin \Art_a\bigr).
\]
The diagonal $\delta(a \to a)$ reproduces papers I--II exactly (a useful
end-to-end validation, since the present computation is an independent
implementation). Off the diagonal we find
(Section~\ref{sec:stats}):

\begin{enumerate}
\item \emph{Sign and size.} $101$ of the $132$ ordered off-diagonal pairs
  have $\delta(a \to b) < 0$, with mean $-0.0035$ --- repulsion persists
  across bases, but at an order of magnitude weaker than on the diagonal
  (mean $-0.034$ over the eleven bases with negative diagonal). The residue channels that transmit the
  Lemke Oliver--Soundararajan bias~\cite{LOS2016} are mostly \emph{different}
  for $\chi_a$ and $\chi_b$, so most of the effect cancels.
\item \emph{Significance.} $124$ of $132$ off-diagonal correlations have
  $|z| > 3$; the structure is small but real.
\item \emph{Asymmetry.} $\delta(a \to b) \neq \delta(b \to a)$: the mean
  absolute difference is $0.009$ and the maximum $0.051$ (attained by
  $(6, 10)$ versus $(10, 6)$). No same-prime analogue of this exists ---
  Paper III's $\phi(a,b)$ is symmetric by construction --- and the asymmetry
  witnesses the ordered conditioning through $p_n - 1$ versus
  $p_{n+1} - 1$.
\item \emph{Extremes are entangled or share structure.} The most negative
  mixed correlations are $\delta(6 \to 2) = -0.0354$ and
  $\delta(21 \to 5) = -0.0330$; the most positive are
  $\delta(6 \to 10) = +0.0415$ and $\delta(3 \to 15) = +0.0314$. All four
  involve multiplicatively dependent pairs. Across all $132$ ordered pairs,
  entanglement of $\{a, b\}$ correlates with $|\delta(a \to b)|$ at
  $r = 0.36$, the strongest single organising variable we find, followed by
  the conductor of the \emph{second} base ($r(\log f_b, |\delta|) = -0.29$);
  the conductor of the first base is essentially irrelevant
  ($r = 0.07$).
\end{enumerate}

\subsection*{A base with positive diagonal correlation}
Extending the base set to include $a = 15$ (absent from papers I--II)
produced a surprise worth reporting on its own: base $15$ has
\[
  \delta(15 \to 15) = +0.0122 \qquad (z = 87.0),
\]
the first known base whose consecutive-prime Artin correlation is globally
\emph{positive}. Papers I--II measured $\delta < 0$ for all eleven bases
examined and conjectured persistence of the negative sign; base $15$ refutes
the natural extension of that conjecture to all bases. The mechanism is the
class-competition phenomenon identified in paper II: base $15$ has three
exclusion classes modulo $60$ (forcing $\delta$ negative on $9.4\%$ of
pairs), but its remaining gap classes are dominated by strongly positively
coupled ones (e.g.\ $\delta = +0.64$ at $g \equiv 10 \pmod{60}$, $+0.26$ at
$g \equiv 18$), and at conductor $60$ --- the largest in our set --- the
negative within-class residue effect is too weak to overcome them. The sign
of $\delta(a)$ is thus not universal: it is the resolution of a competition
that either side can win, and $15$ is the smallest base where the positive
side wins.

\section{The mixed exclusion classification}
\label{sec:theory}

Throughout, $d_a = \sqf(a)$, $f_a$ is the conductor of
$\chi_a = \bigl(\tfrac{d_a}{\cdot}\bigr)$ (namely $d_a$ if
$d_a \equiv 1 \bmod 4$, else $4 d_a$), and $L = \mathrm{lcm}(f_a, f_b)$.
Admissibility of $r$ means $\gcd(r, L) = 1$; residues attained by primes
$p \nmid 2ab$ are exactly the admissible ones.

\begin{definition}
\label{def:mixed}
The class $g \bmod L$ is a \emph{mixed exclusion class} for the ordered pair
$(a, b)$ if no admissible $r$ with $r + g$ admissible has
$\chi_a(r) = -1$ and $\chi_b(r + g) = -1$.
\end{definition}

\begin{proposition}
\label{prop:mixedlaw}
If $g$ is a mixed exclusion class for $(a, b)$, then there are no primes
$p, q = p + g'$ with $p \nmid 2ab$, $q \nmid 2ab$, $g' \equiv g \pmod L$,
such that $a$ is a primitive root modulo $p$ and $b$ is a primitive root
modulo $q$.
\end{proposition}

\begin{proof}
$\chi_a(p)$ depends only on $p \bmod f_a$, hence on $p \bmod L$; similarly
for $\chi_b(q)$. If both primitive-root conditions held, then
$\chi_a(p) = \chi_b(q) = -1$ with $p \bmod L$ admissible and
$q \equiv p + g$, contradicting Definition~\ref{def:mixed}.
\end{proof}

The classification of mixed exclusion classes reduces, exactly as in paper
II, to the behaviour of the prime-discriminant components of $\chi_a$ and
$\chi_b$ under the shift, together with the inadmissibility counting at odd
prime conductors; we do not repeat the component analysis, but record the
outcome of the exhaustive scan for our twelve bases
$\{2, 3, 5, 6, 7, 10, 11, 13, 15, 17, 21, 29\}$: the ordered pairs
possessing mixed exclusion classes are the twelve diagonal pairs (with the
classes found in papers I--II, including the base-$5$ inadmissibility classes
and both classes of base $21$), together with exactly one unordered
off-diagonal pair:

\begin{theorem}
\label{thm:mixed26}
The ordered pairs $(2, 6)$ and $(6, 2)$ have the mixed exclusion classes
$g \equiv 10, 14 \pmod{24}$; no other off-diagonal pair from the twelve bases
admits any mixed exclusion class.
\end{theorem}

\begin{proof}[Proof for $(2,6)$; the converse half is the exhaustive scan]
$\chi_2$ has conductor $8$ with $\chi_2(r) = -1$ exactly for
$r \equiv 3, 5 \pmod 8$; $\chi_6$ has conductor $24$ with
$\chi_6(r) = -1$ exactly for $r \equiv 7, 11, 13, 17 \pmod{24}$. Fix
$g \equiv 10 \pmod{24}$. The residues $r \bmod 24$ with $\chi_2(r) = -1$ are
$r \equiv 3, 5, 11, 13, 19, 21$; discarding those not coprime to $24$ leaves
$r \equiv 5, 11, 13, 19$. Their shifts $r + 10$ are
$15, 21, 23, 5 \pmod{24}$; the first two are inadmissible ($3 \mid 15, 21$),
and $\chi_6(23) = \chi_6(5) = +1$. So no admissible configuration has both
characters $-1$. The case $g \equiv 14$ is symmetric ($r + 14 \equiv
19, 1, 3, 9$; admissible ones are $19, 1$, and $\chi_6(19) = \chi_6(1) =
+1$). Note the essential role of inadmissibility: on the full residue system
the configurations would exist, but they land on multiples of $3$.
\end{proof}

\begin{remark}
The pair $(2, 6)$ is multiplicatively dependent ($\sqf(2 \cdot 6) = 3$,
present in the ambient set), and the mechanism visibly mixes the two
ingredients of paper II: a character shift (as in reversal) and residue
inadmissibility (as at base $5$). That the same two mechanisms exhaust the
mixed case as well --- confirmed by the converse scan below --- supports the
view of papers II and III that quadratic-character parity and admissibility
are the only deterministic constraints at this level, with all further
structure being statistical.
\end{remark}

\subsection*{Verification at $10^9$ and completeness}
For every ordered pair and every predicted class we counted, among the
$50{,}847{,}531$ consecutive prime pairs below $10^9$, the pairs whose gap
lies in the class and, among them, those with $a$ a primitive root of $p_n$
and $b$ of $p_{n+1}$. Table~\ref{tab:verification} reports the nonzero rows.
The predicted classes contain $104{,}972{,}810$ pairs in total and the
co-occurrence count is zero in every row. Conversely, a data-driven scan of
\emph{all} classes of \emph{all} $144$ ordered pairs (modulo the appropriate
$L$) for zero counts supported by at least $5 \times 10^5$ pairs returns
exactly the predicted classes and nothing else: the classification is
complete at the resolution of the data.

\begin{table}[t]
\centering
\caption{Verification of mixed exclusion classes at $x = 10^9$. Ordered pairs
$(a, b)$: $a$ tested at $p_n$, $b$ at $p_{n+1}$. Diagonal rows reproduce
papers I--II (including the inadmissibility classes of base $5$); the
$(2,6)/(6,2)$ rows are the genuinely mixed law of Theorem~\ref{thm:mixed26}.}
\label{tab:verification}
\begin{tabular}{llrrr}
\toprule
$(a, b)$ & classes mod $L$ & $L$ & pairs & co-occurrences \\
\midrule
$(2,2)$   & $4$          & $8$  & $13{,}404{,}106$ & $0$ \\
$(3,3)$   & $4, 6, 8$    & $12$ & $27{,}010{,}759$ & $0$ \\
$(5,5)$   & $2, 3$       & $5$  & $20{,}523{,}554$ & $0$ \\
$(6,6)$   & $8, 12, 16$  & $24$ & $12{,}419{,}039$ & $0$ \\
$(7,7)$   & $14$         & $28$ & $3{,}567{,}335$  & $0$ \\
$(10,10)$ & $20$         & $40$ & $2{,}214{,}511$  & $0$ \\
$(11,11)$ & $22$         & $44$ & $1{,}796{,}191$  & $0$ \\
$(15,15)$ & $20, 30, 40$ & $60$ & $4{,}782{,}744$  & $0$ \\
$(21,21)$ & $7, 14$      & $21$ & $4{,}046{,}375$  & $0$ \\
$(2,6)$   & $10, 14$     & $24$ & $7{,}604{,}098$  & $0$ \\
$(6,2)$   & $10, 14$     & $24$ & $7{,}604{,}098$  & $0$ \\
\midrule
\emph{total} & & & $104{,}972{,}810$ & $\mathbf{0}$ \\
\bottomrule
\end{tabular}
\end{table}

\section{Data and computation}
\label{sec:data}

The computation is one streaming pass of a segmented sieve to $10^9$
($50{,}847{,}532$ primes $\ge 5$; $50{,}847{,}531$ consecutive pairs). For
each prime, $p - 1$ is factored once by trial division and the Artin
indicator computed for all twelve bases by the standard criterion
($a^{(p-1)/\ell} \not\equiv 1$ for all prime $\ell \mid p-1$, $128$-bit
arithmetic); for each ordered base pair $(a, b)$ we accumulate the
$2 \times 2$ table of $\bigl(\mathbf{1}[p_n \in \Art_a],
\mathbf{1}[p_{n+1} \in \Art_b]\bigr)$, globally and per gap $g \le 300$. All
counts are exact. Hardware and validation are as in the companion papers
($2$-core Xeon E5-2673~v4, $8$~GB RAM; single core); the run takes under an
hour. Beyond the usual checks (prime counts; marginal densities matching the
Hooley values to four decimals, including $0.393642$ for base $5$ against
$C \cdot 20/19 = 0.393638$), the twelve diagonal correlations reproduce the
values of papers I--II to all reported digits from an independently written
accumulator, and the verification of Table~\ref{tab:verification} doubles as
a check of the classification.

\section{The mixed correlation matrix}
\label{sec:stats}

\begin{table}[t]
\centering
\caption{Extremes of the off-diagonal mixed correlation
$\delta(a \to b)$ at $x = 10^9$. ``Dep.''\ marks multiplicatively dependent
$\{a, b\}$ (the product is a square times a third base of the set).}
\label{tab:extremes}
\begin{tabular}{rrrrc}
\toprule
$a$ & $b$ & $\delta(a \to b)$ & $z$ & Dep. \\
\midrule
$6$  & $2$  & $-0.0354$ & $-252.6$ & \checkmark \\
$21$ & $5$  & $-0.0330$ & $-233.1$ & \checkmark \\
$3$  & $7$  & $-0.0233$ & $-165$ & \checkmark \\
$11$ & $3$  & $-0.0201$ & $-142$ & \\
$7$  & $3$  & $-0.0193$ & $-137$ & \checkmark \\
\midrule
$2$  & $6$  & $+0.0094$ & $+67.2$ & \checkmark \\
$15$ & $3$  & $+0.0222$ & $+158.4$ & \checkmark \\
$3$  & $15$ & $+0.0314$ & $+224.2$ & \checkmark \\
$6$  & $10$ & $+0.0415$ & $+295.8$ & \checkmark \\
\bottomrule
\end{tabular}
\end{table}

The full $12 \times 12$ matrix is available with the data; we summarise its
structure.

\begin{observation}[Weak, mostly negative, highly significant]
\label{obs:offdiag}
Of the $132$ ordered off-diagonal pairs, $101$ have $\delta(a \to b) < 0$;
the mean is $-0.0035$, one order of magnitude below the mean $-0.034$ of
the eleven negative diagonal values; and $124$ of $132$ have $|z| > 3$. As in the companion papers, at
this sample size the $z$-scores certify only that the values are real; the
relevant magnitudes are the effect sizes.
\end{observation}

The interpretation follows the channel picture of papers I--II: the
Lemke Oliver--Soundararajan repulsion acts on the joint residues of
$(p_n, p_{n+1})$, and transmits to the pair
$(\Art_a(p_n), \Art_b(p_{n+1}))$ only through residue classes that both
characters (and both $\omega$-channels) see. For $a = b$ every channel is
shared; for $a \neq b$ only the common structure --- the mutual conductor
divisors and the universal $\omega(p-1)$ channel --- survives, and the
correlation drops by the factor observed.

\begin{observation}[Asymmetry]
\label{obs:asym}
$\delta(a \to b) \neq \delta(b \to a)$ in general: the mean of
$|\delta(a \to b) - \delta(b \to a)|$ over unordered pairs is $0.0090$ and
the maximum is $0.0511$, attained by $\{6, 10\}$
($\delta(6 \to 10) = +0.0415$ versus $\delta(10 \to 6) = -0.0096$).
\end{observation}

The asymmetry has no analogue in paper III (where the same-prime $\phi(a,b)$
is symmetric by definition) and directly witnesses the ordering of the
conditioning: $\delta(a \to b)$ conditions on the arithmetic of $p_n - 1$
and measures on $p_{n+1} - 1$. Its size --- comparable to the off-diagonal
effects themselves --- shows that the mixed correlation is not a symmetric
``base similarity'' but a directed transition statistic.

\begin{observation}[What organises the magnitudes]
\label{obs:organise}
Over the $132$ ordered pairs: multiplicative dependence of $\{a, b\}$
correlates with $|\delta(a \to b)|$ at $r = +0.36$; the conductor of the
second base at $r(\log f_b, |\delta|) = -0.29$ (and
$r(\log \mathrm{lcm}(f_a, f_b), |\delta|) = -0.32$); the conductor of the
first base is negligible ($r = +0.07$). All four extreme values of
Table~\ref{tab:extremes} involve dependent pairs.
\end{observation}

That $f_b$ matters and $f_a$ does not is consistent with the conductor law
of paper II acting at the \emph{measured} prime: the size of
$|\delta(a \to b)|$ is controlled by how concentrated the $\chi_b$-channels
are at $p_{n+1}$, while the conditioning event at $p_n$ enters mainly
through whether its channels overlap those of $b$ --- which is what
multiplicative dependence measures.

\section{Base 15: a positive diagonal}
\label{sec:base15}

Base $15$ was not among the eleven bases of papers I--II. Its diagonal
correlation at $x = 10^9$ is
\[
  \delta(15 \to 15) = +0.012196, \qquad z = +87.0,
\]
positive --- the first base known to us with globally positive
consecutive-prime Artin correlation, against the uniform negativity of the
other twelve measured so far (this paper adds $15$; papers I--II measured
eleven others, all negative).

The decomposition by gap class modulo $60$ (conductor of $\chi_{15}$)
explains the sign. The exclusion classes $g \equiv 20, 30, 40 \pmod{60}$
contain $9.4\%$ of all pairs and force $P(\text{both Artin}) = 0$ there; but
among the populous non-excluded classes the coupling is strongly positive:
$\delta = +0.64$ at $g \equiv 10$, $+0.26$ at $g \equiv 18$, $+0.24$ at
$g \equiv 6, 8$, against $-0.20$ at $g \equiv 2$ and $-0.19$ at
$g \equiv 12$. In the class-competition language of paper II, base $15$ is
the first base whose positively coupled classes outweigh the forced zeros
and the within-class repulsion. Two structural facts make this possible: the
conductor $60$ is the largest in the set (so the within-class
Lemke Oliver--Soundararajan residue effect, which scales down with
conductor, is the weakest), while the character $\chi_{15} = \chi_{-3}
\chi_{5} \chi_{-4}$ has three prime-discriminant components whose joint
preservation classes ($g \equiv 0, 10, 50 \pmod{60}$, plus partial
preservation elsewhere) are unusually rich. We conjecture that
$\delta(a) > 0$ for infinitely many bases --- specifically for bases whose
conductor is large and whose preserving classes are dense --- and note that
the sign of $\delta(a)$ is therefore not an intrinsic feature of the Artin
property but an arithmetic accident of the base.

\section{Discussion}
\label{sec:discussion}

\subsection*{The programme completed}
Across the four papers the picture is uniform. On every axis the joint
behaviour of Artin status splits into a deterministic part --- parity and
admissibility constraints on quadratic characters, now classified in all
three non-trivial settings (diagonal consecutive, same-prime cross-base,
mixed) --- and a statistical part carried by the arithmetic of $p - 1$: the
LOS bias through residue channels for consecutive primes, the shared
factorisation for same-prime pairs, and their directed, mostly cancelling
combination in the mixed case. The deterministic laws are absolute but
concentrated on sparse gap classes; the statistical effects are small but
global; and which of them dominates a headline number like the sign of
$\delta(a)$ is a competition, as base $15$ now demonstrates.

\subsection*{Future work}
(1) The channel-by-channel derivation of the conductor law (paper II)
extends naturally to the mixed setting and would predict the $f_b$-dominance
of Observation~\ref{obs:organise}; the mixed data provide $132$ constraints
for such a derivation rather than eleven. (2) A systematic search for further
positive-$\delta$ bases (candidates: large conductor, many prime-discriminant
components --- e.g.\ $a = 30, 35, 39, 105$) would test the conjecture of
Section~\ref{sec:base15}. (3) The decay in $x$ of the mixed correlations,
and of $\delta(15)$ in particular, remains unmeasured; the sign of
$\delta(15)$ at $10^{12}$ is an especially clean question, since the
LOS-driven within-class repulsion decays like $1/\log x$ while the
class-composition effect does not.

\section*{Acknowledgements}
The author thanks the developers of the open-source tools used in this work
(GCC, Python, \texttt{sympy}, \texttt{matplotlib}).

\subsection*{Use of generative AI}
The author used Anthropic's Claude (models Claude Opus 4.5 and Claude
Opus 5) as an assistant: for drafting and revising exposition, for writing
and debugging the sieve and analysis programs, and for exploratory
discussion. The mixed exclusion classification was generated by an
exhaustive symbolic scan and verified both by hand (Theorem~\ref{thm:mixed26})
and against the full dataset; it was this scan, run in preparation for the
present paper, that uncovered the incompleteness of the original
classification of paper II (since corrected there and acknowledged in that
paper). All computations were executed and verified on the author's own
hardware; every reference has been verified against the published record.
The author is solely responsible for the content.

\section*{Declaration of interest}
The author reports there are no competing interests to declare.

\section*{Funding}
This research received no external funding.

\section*{Data availability}
All code and results are publicly available at
\url{https://github.com/jpbald93/consecutive-artin}: the one-pass program
(\texttt{mixed\_axis.c}), the symbolic classification scripts, the full
$12 \times 12$ correlation matrix and per-gap contingency tables, and
scripts reproducing every number in this paper. The computation reproduces
in under an hour on a single core.

\begin{thebibliography}{9}

\bibitem{BaldI}
J.~Bald,
\emph{Correlations between primitive root statuses of consecutive primes},
preprint (2026); code and data at
\url{https://github.com/jpbald93/consecutive-artin}.

\bibitem{BaldII}
J.~Bald,
\emph{Exclusion laws for consecutive Artin primes in arbitrary bases},
preprint (2026); code and data at
\url{https://github.com/jpbald93/consecutive-artin}.

\bibitem{BaldIII}
J.~Bald,
\emph{Cross-base correlations of Artin primes: entanglement, exclusion, and
the structure of $p-1$},
preprint (2026); code and data at
\url{https://github.com/jpbald93/consecutive-artin}.

\bibitem{Hooley1967}
C.~Hooley,
\emph{On Artin's conjecture},
J. Reine Angew. Math. \textbf{225} (1967), 209--220.

\bibitem{LOS2016}
R.~J. Lemke~Oliver and K.~Soundararajan,
\emph{Unexpected biases in the distribution of consecutive primes},
Proc. Natl. Acad. Sci. USA \textbf{113} (2016), no.~31, E4446--E4454.

\end{thebibliography}

\end{document}
